\section*{Capítulo \ref{ch:g9:natural-risks} --- Vivir con los riesgos naturales}

\begin{solution}{exo:g9:natural-risks:1}
(a) Peligro. (b) \omterm{def:g9:natural-risks:vulnerability}{Vulnerabilidad}. (c) \omterm{def:g9:natural-risks:vulnerability}{Vulnerabilidad}. (d) Peligro. (e) Las
dos cosas: las coladas de \omterm{def:g3:stones-pebbles-sand:sand}{barro} son el peligro; construir la ciudad en su
camino la hace vulnerable.
\end{solution}

\begin{solution}{exo:g9:natural-risks:2}
(a) Falso: donde no vive nadie, nada puede sufrir daños, así que el \omterm{def:g9:natural-risks:risk}{riesgo}
es bajo. (b) Verdadero. (c) Verdadero. (d) Falso: una alerta no cambia el
\omterm{def:g4:volcanoes-earthquakes-first-look:volcano}{volcán}; reduce la \omterm{def:g9:natural-risks:vulnerability}{vulnerabilidad} de la población, que puede marcharse a
tiempo.
\end{solution}

\begin{solution}{exo:g9:natural-risks:3}
(a) Sí, con días de antelación. (b) Sí, con horas o días de antelación.
(c) No: los \omterm{def:g4:volcanoes-earthquakes-first-look:earthquake}{terremotos} no se pueden predecir. (d) Sí: una vez que el
\omterm{def:g7:earthquakes:tsunami}{tsunami} se ha originado, tarda horas en cruzar el \omterm{def:g2:water-on-earth:ocean}{océano}, y la alerta
viaja más deprisa.
\end{solution}

\begin{solution}{exo:g9:natural-risks:4}
(a) Las sacudidas del suelo: enjambres de pequeños \omterm{def:g4:volcanoes-earthquakes-first-look:earthquake}{terremotos} cuando el
\omterm{def:g7:volcanoes:magma}{magma} sube. (b) La inclinación de la ladera: el \omterm{def:g4:volcanoes-earthquakes-first-look:volcano}{volcán} que se hincha.
(c) Cómo se desplaza el suelo, milímetro a milímetro: la hinchazón a medida
que se acumula el \omterm{def:g7:volcanoes:magma}{magma}. (d) Los gases del penacho: más gas que escapa
cuando el \omterm{def:g7:volcanoes:magma}{magma} asciende.
\end{solution}

\begin{solution}{exo:g9:natural-risks:5}
(a) $45 \div 2 = \qty{22.5}{km/h}$. (b) $4 \times 2 = \qty{8}{km}$,
suficiente para llegar a un terreno más alto. (c) Nadie vio la \omterm{def:g4:volcanoes-earthquakes-first-look:eruption}{erupción},
oculta por la \omterm{def:g5:natural-hazards:storm}{tormenta}, y nadie los avisó: no había ningún medio rápido de
hacer llegar la información a la ciudad, ni plan de evacuación. (d) Por
ejemplo: un plan de evacuación con rutas hacia terrenos más altos,
practicado de antemano; una cadena de alerta de los científicos a la
ciudad; más instrumentos en el \omterm{def:g4:volcanoes-earthquakes-first-look:volcano}{volcán}.
\end{solution}

\begin{solution}{exo:g9:natural-risks:6}
(a) La \omterm{def:g5:water-shapes-the-land:flood}{crecida} de \qty{10}{m} sobre el \omterm{def:g1:landscapes-around-us:water}{río}. (b) Unos \qty{0.47}{km} para
la \omterm{def:g5:water-shapes-the-land:flood}{crecida} de \qty{2}{m}, \qty{0.56}{km} para la de \qty{5}{m} y
\qty{0.67}{km} para la de \qty{10}{m} (son válidas las lecturas con un
margen de unos \qty{0.05}{km}). (c) Solo la \omterm{def:g5:water-shapes-the-land:flood}{crecida} de \qty{10}{m}: el
terreno está a más de 2 y de \qty{5}{m} sobre el \omterm{def:g1:landscapes-around-us:water}{río}. (d) Basta una
pequeña subida del \omterm{def:g1:landscapes-around-us:water}{río} para alcanzarla, y las \omterm{def:g5:water-shapes-the-land:flood}{crecidas} pequeñas llegan
más a menudo que las grandes.
\end{solution}

\begin{solution}{exo:g9:natural-risks:7}
(a) Peligro alto, \omterm{def:g9:natural-risks:vulnerability}{vulnerabilidad} baja: \omterm{def:g9:natural-risks:risk}{riesgo} moderado. (b) Peligro bajo,
\omterm{def:g9:natural-risks:vulnerability}{vulnerabilidad} alta: \omterm{def:g9:natural-risks:risk}{riesgo} moderado. (c) Peligro alto, \omterm{def:g9:natural-risks:vulnerability}{vulnerabilidad}
alta: \omterm{def:g9:natural-risks:risk}{riesgo} alto. (d) Peligro bajo, \omterm{def:g9:natural-risks:vulnerability}{vulnerabilidad} baja: \omterm{def:g9:natural-risks:risk}{riesgo} bajo.
\end{solution}

\begin{solution}{exo:g9:natural-risks:8}
(a) C, después B, después A. (b) Las riostras forman triángulos, que no
pueden inclinarse sin estirar o aplastar una riostra: la estructura
mantiene su forma en lugar de inclinarse. (c) Los aisladores ceden cuando
el suelo se mueve de un lado a otro bajo ellos, de modo que a la escuela le
llega mucho menos movimiento.
\end{solution}

\begin{solution}{exo:g9:natural-risks:9}
(a) $1 \div 100 = 1~\%$ al año. (b) La probabilidad es la misma cada año,
pase lo que pase el año anterior: la \omterm{def:g5:water-shapes-the-land:flood}{crecida} puede volver el año que
viene. (c) $300 \div 100 = 3$ \omterm{def:g5:water-shapes-the-land:flood}{crecidas}, por término medio.
\end{solution}

\begin{solution}{exo:g9:natural-risks:10}
(a) $3600 \div 15 = 240$ lecturas. (b) $240 \times 24 \times 365 =
\num{2102400} \approx 2.1 \times 10^6$ lecturas. (c) $2400 \div 800 =
\qty{3}{h}$. (d) En esas tres horas, los centros de alerta pueden avisar a
la \omterm{def:g1:landscapes-around-us:water}{costa} y la población puede ponerse a salvo en terrenos más altos.
\end{solution}

\begin{solution}{exo:g9:natural-risks:11}
(a) $4 \times 3 \times 3.8 = \qty{45.6}{L}$. (b) $45.6 \div 1.5 = 30.4$,
así que 31 botellas. (c) Por ejemplo: comida que se conserve, para comer
cuando las tiendas estén cerradas; una linterna, por si se corta la
electricidad; una radio de pilas o de manivela, para oír las alertas; un
botiquín, para curar heridas leves; mantas de abrigo. (d) Para poder
llevárselo enseguida si la familia tiene que marcharse deprisa.
\end{solution}

\begin{solution}{exo:g9:natural-risks:12}
Por ejemplo: (1) elaborar un \omterm{def:g9:natural-risks:hazard-map}{mapa de peligrosidad} de las coladas de \omterm{def:g3:stones-pebbles-sand:sand}{barro},
\omterm{def:g1:landscapes-around-us:landform}{valle} por \omterm{def:g1:landscapes-around-us:landform}{valle} (lado del peligro); (2) instalar \omterm{def:g7:earthquakes:seismometer}{sismómetros},
inclinómetros, receptores GPS e instrumentos de gases en el \omterm{def:g4:volcanoes-earthquakes-first-look:volcano}{volcán} (lado
del peligro); (3) establecer una cadena de alerta rápida del observatorio
a la ciudad, con sirenas (lado de la \omterm{def:g9:natural-risks:vulnerability}{vulnerabilidad}); (4) planificar y
señalizar rutas hacia terrenos más altos, y practicarlas con simulacros
(lado de la \omterm{def:g9:natural-risks:vulnerability}{vulnerabilidad}); (5) no construir nuevas viviendas ni escuelas
en el camino de las coladas de \omterm{def:g3:stones-pebbles-sand:sand}{barro} (lado de la \omterm{def:g9:natural-risks:vulnerability}{vulnerabilidad}).
\end{solution}

\begin{solution}{pb:g9:natural-risks:1}
\textbf{1.} El \omterm{def:g1:landscapes-around-us:water}{río} está a \qty{143}{m} junto a C, a \qty{137}{m} junto a
A y a \qty{126}{m} junto a B; en el \omterm{def:g2:drawing-to-map:map}{mapa}, C está al oeste y B al este, así
que el \omterm{def:g1:landscapes-around-us:water}{río} fluye hacia el este.

\textbf{2.} A: $137 - 137 = \qty{0}{m}$; B: $276 - 126 = \qty{150}{m}$;
C: $236 - 143 = \qty{93}{m}$.

\textbf{3.} Solo A, que está a menos de \qty{5}{m} sobre el \omterm{def:g1:landscapes-around-us:water}{río}.

\textbf{4.} Unos \qty{0.57}{km} (entre 0.5 y \qty{0.65}{km}).

\textbf{5.} La elevación en \qty{100}{m} es la mitad de la elevación en
\qty{200}{m}: A $1.4 \div 2 = 0.7~\%$; B $83 \div 2 = 41.5~\%$;
C $20 \div 2 = 10~\%$.

\textbf{6.} B: su pendiente es muy fuerte, y la \omterm{def:g3:stones-pebbles-sand:rock}{roca} suelta y el suelo
pueden deslizarse por una pendiente fuerte, sobre todo tras \omterm{def:g4:clouds-rain-snow:precipitation}{lluvias}
intensas.

\textbf{7.} A: alrededor de \qty{1}{km}; B: casi sobre la \omterm{def:g7:earthquakes:fault}{falla} (a menos
de \qty{0.05}{km}); C: alrededor de \qty{1.3}{km}.

\textbf{8.} Si la \omterm{def:g7:earthquakes:fault}{falla} se mueve, el suelo se rompe y se desplaza a lo
largo de ella, lo que puede partir un edificio; y un \omterm{def:g4:volcanoes-earthquakes-first-look:earthquake}{terremoto} en la \omterm{def:g7:earthquakes:fault}{falla}
sacudiría allí el suelo con mucha fuerza.

\textbf{9.}

\begin{center}
{\small
\begin{tabular}{|c|c|c|c|}
\hline
emplazamiento & ¿en la \omterm{def:g5:water-shapes-the-land:flood}{zona inundable}? & ¿pendiente fuerte? & ¿sobre la \omterm{def:g7:earthquakes:fault}{falla}? \\
\hline
A & sí & no & no \\
B & no & sí & sí \\
C & no & no & no \\
\hline
\end{tabular}}
\end{center}

\textbf{10.} El emplazamiento C: es el único con tres ``no'': \qty{93}{m}
sobre el \omterm{def:g1:landscapes-around-us:water}{río}, una pendiente suave del \qty{10}{\%} y a más de \qty{1}{km}
de la \omterm{def:g7:earthquakes:fault}{falla}.

\textbf{11.} Por ejemplo: riostras diagonales y muros de cortante en la
estructura, uniones sólidas, o aisladores bajo el edificio; los alumnos
practican simulacros de \omterm{def:g4:volcanoes-earthquakes-first-look:earthquake}{terremoto} (agáchate, cúbrete, agárrate).

\textbf{12.} \textbf{17 casas} de 28 están dentro de la \omterm{def:g5:water-shapes-the-land:flood}{zona inundable}:
$17 \div 28 \approx 0.61$, unas 3 casas de cada 5.
\end{solution}
